%! Lengths and Angles from Dot Products — Unit 1, Lesson 1.2 — Guided Notes
\documentclass[10pt]{article}
\usepackage{linalg-article}
\usepackage{linalg-boxes}
\newcommand{\TallMath}[1]{$\displaystyle #1\rule[-1.4em]{0pt}{3.2em}$}
\newcommand{\vv}{\mathbf{v}}
\newcommand{\ww}{\mathbf{w}}

\begin{document}

\pageheader{Unit 1, Lesson 1.2}{Guided Notes}
\namedateperiod

\vspace{0.12in}

\begin{objectivebox}
\textbf{By the end of this lesson, I will be able to\dots}
\begin{itemize}\setlength{\itemsep}{2pt}
  \item compute the \emph{dot product} $\vv\cdot\ww$ and use it to find a vector's \emph{length};
  \item find the \emph{angle} between two vectors with the cosine formula;
  \item decide when two vectors are \emph{perpendicular} and read the sign of the dot product.
\end{itemize}
\end{objectivebox}

\vspace{0.10in}

\begin{vocabbox}
\small
Fill in each term as we build it together.
\vspace{2pt}
\termblanklong{Dot product}
\termblanklong{Length (norm)}
\termblanklong{Unit vector}
\termblanklong{Angle formula}
\termblanklong{Perpendicular (orthogonal)}
\end{vocabbox}

\begin{hookbox}
Two pairs of arrows start at the origin. One pair points nearly the \emph{same} way; the other
pair meets at a \emph{right angle}.
\begin{itemize}\setlength{\itemsep}{1pt}
  \item Pair A: $\vv=\begin{bsmallmatrix} 3\\4 \end{bsmallmatrix}$ and $\ww=\begin{bsmallmatrix} 4\\3 \end{bsmallmatrix}$.
        Multiply matching components and add: $(3)(4)+(4)(3)=\blank{2.0cm}$
  \item Pair B: $\begin{bsmallmatrix} 3\\1 \end{bsmallmatrix}$ and $\begin{bsmallmatrix} -1\\3 \end{bsmallmatrix}$.
        Do the same: $(3)(-1)+(1)(3)=\blank{2.0cm}$
\end{itemize}
Which pair gave the bigger number? Which gave zero? \writeline
\end{hookbox}

\begin{notesbox}{1. What Is the Dot Product?}
The \textbf{dot product} of $\vv=\begin{bmatrix} v_1 \\ v_2 \end{bmatrix}$ and
$\ww=\begin{bmatrix} w_1 \\ w_2 \end{bmatrix}$ multiplies \emph{matching} components and adds:
\[
  \vv\cdot\ww = v_1 w_1 + v_2 w_2 .
\]
The answer is a single \textbf{number} (a scalar) --- \emph{not} a vector.

\vspace{2pt}
\textit{Example.} With $\vv=\begin{bmatrix} 3 \\ 4 \end{bmatrix}$ and $\ww=\begin{bmatrix} 4 \\ 3 \end{bmatrix}$:
\[
  \vv\cdot\ww = (3)(4) + (4)(3) = \blank{2.4cm}.
\]

\vspace{2pt}
Your turn: $\begin{bmatrix} 2 \\ -1 \end{bmatrix}\cdot\begin{bmatrix} 3 \\ 5 \end{bmatrix}
= (2)(\underline{\hspace{0.8cm}}) + (-1)(\underline{\hspace{0.8cm}}) = \blank{2.0cm}$.
\end{notesbox}

\begin{notesbox}{2. Length Is a Dot Product with Itself}
Dot a vector with \emph{itself} and every term becomes a square:
\[
  \vv\cdot\vv = v_1^2 + v_2^2 , \qquad\text{so}\qquad \|\vv\| = \sqrt{\vv\cdot\vv}.
\]
This is exactly the Lesson~1.0 length formula, now written with a dot product.

\vspace{2pt}
\textit{Example.} For $\vv=\begin{bmatrix} 3 \\ 4 \end{bmatrix}$:
$\vv\cdot\vv = 9 + 16 = \blank{1.4cm}$, so $\|\vv\| = \sqrt{\,\blank{1.2cm}\,} = \blank{1.2cm}$.

\vspace{2pt}
A \textbf{unit vector} has length $1$. To make one, scale $\vv$ by $\dfrac{1}{\|\vv\|}$:
the unit vector along $\begin{bsmallmatrix} 3\\4 \end{bsmallmatrix}$ is
$\dfrac{1}{\blank{0.8cm}}\begin{bmatrix} 3 \\ 4 \end{bmatrix}
= \begin{bmatrix}\rule{0pt}{1.1em}\blank{1.0cm}\\[2pt]\blank{1.0cm}\end{bmatrix}$.
\end{notesbox}

\begin{notesbox}{3. The Dot Product Finds the Angle}
The dot product hides the angle $\theta$ between the two vectors:
\[
  \vv\cdot\ww = \|\vv\|\,\|\ww\|\cos\theta
  \qquad\Longrightarrow\qquad
  \cos\theta = \frac{\vv\cdot\ww}{\|\vv\|\,\|\ww\|}.
\]
\begin{center}
\begin{tikzpicture}[scale=0.7]
  \draw[->, thick, charcoal] (-0.4,0)--(3.4,0) node[right]{\scriptsize $x$};
  \draw[->, thick, charcoal] (0,-0.4)--(0,3.4) node[above]{\scriptsize $y$};
  \draw[->, very thick, burgundy] (0,0)--(3,0) node[midway, below]{$\vv$};
  \draw[->, very thick, royalblue] (0,0)--(2,2) node[above right]{$\ww$};
  \draw[burgundy!70] (0.7,0) arc (0:45:0.7);
  \node at (1.15,0.35) {\scriptsize $\theta$};
\end{tikzpicture}
\end{center}
\textit{Example.} $\vv=\begin{bmatrix} 2 \\ 0 \end{bmatrix}$, $\ww=\begin{bmatrix} 2 \\ 2 \end{bmatrix}$:
$\vv\cdot\ww=4$, $\|\vv\|=2$, $\|\ww\|=2\sqrt2$, so
$\cos\theta=\dfrac{4}{(2)(2\sqrt2)}=\dfrac{1}{\sqrt2}$ and $\theta=\blank{1.6cm}$.

\vspace{4pt}
\textbf{The sign tells the story.} Fill in the angle each sign forces:
\begin{center}\small
\renewcommand{\arraystretch}{1.2}
\begin{tabular}{c|c|c}
$\vv\cdot\ww > 0$ & $\vv\cdot\ww = 0$ & $\vv\cdot\ww < 0$ \\ \hline
angle is \blank{1.6cm} & angle is \blank{1.6cm} & angle is \blank{1.6cm} \\
(same-ish way) & (right angle) & (opposing) \\
\end{tabular}
\end{center}
\end{notesbox}

\begin{notesbox}{4. Perpendicular Means Dot Product Zero}
Since $\cos 90^\circ = 0$, the angle formula makes one case special: two vectors are
\textbf{perpendicular} (meet at a right angle) exactly when
\[
  \vv\cdot\ww = 0 .
\]
\begin{center}
\begin{tikzpicture}[scale=0.6]
  \draw[step=1, linegray!60, thin] (-1.4,-0.4) grid (3.4,3.4);
  \draw[->, thick, charcoal] (-1.6,0)--(3.6,0) node[right]{\scriptsize $x$};
  \draw[->, thick, charcoal] (0,-0.6)--(0,3.6) node[above]{\scriptsize $y$};
  \draw[->, very thick, burgundy] (0,0)--(3,1) node[right]{$\begin{bsmallmatrix} 3\\1 \end{bsmallmatrix}$};
  \draw[->, very thick, royalblue] (0,0)--(-1,3) node[above left]{$\begin{bsmallmatrix} -1\\3 \end{bsmallmatrix}$};
  \draw[charcoal] (0.28,0.09) -- (0.19,0.38) -- (-0.09,0.28);
\end{tikzpicture}
\end{center}
\textit{Check.} $\begin{bmatrix} 3 \\ 1 \end{bmatrix}\cdot\begin{bmatrix} -1 \\ 3 \end{bmatrix}
= (3)(-1)+(1)(3) = \blank{1.6cm}$, so the two vectors are \blank{2.8cm}.

\vspace{2pt}
Your turn: Is $\begin{bmatrix} 4 \\ 2 \end{bmatrix}$ perpendicular to $\begin{bmatrix} 1 \\ -2 \end{bmatrix}$?
Compute the dot product and decide. \writeline
\end{notesbox}

\begin{practicebox}
Let $\vv = \begin{bmatrix} 6 \\ 8 \end{bmatrix}$ and $\ww = \begin{bmatrix} 1 \\ 2 \end{bmatrix}$.
\begin{enumerate}[leftmargin=1.6em]
  \item Compute $\vv\cdot\ww$. \hfill \blank{4.0cm}
  \item Find $\|\vv\|$ using $\|\vv\|=\sqrt{\vv\cdot\vv}$. \hfill \blank{4.0cm}
  \item Is $\begin{bmatrix} 2 \\ 3 \end{bmatrix}$ perpendicular to $\begin{bmatrix} 3 \\ -2 \end{bmatrix}$?
        Show how you know. \writeline
\end{enumerate}
\end{practicebox}

\end{document}
